\subsection{Detailed voxelized head-and-neck phantom construction}
A synthetic voxelized head-and-neck phantom was constructed algorithmically to provide a controllable but anatomically structured environment for spatially fractionated radiotherapy (SFRT) planning. The phantom was discretized on a $133 \times 173 \times 120$ Cartesian grid with isotropic $1.5$~mm voxels, corresponding to a physical extent of $20 \times 26 \times 18$~cm. The external contour was formed as the union of ellipsoidal head, neck, shoulder, and upper-thoracic components. Internal anatomy was then added using combinations of ellipsoids, capped cylinders, and polyline-based tubular primitives. These included skull, mandible, maxilla, cervical vertebrae, oral-airway-tracheal structures, esophagus, bilateral parotid and submandibular glands, spinal cord, brain, brainstem, and a thin blood-brain barrier (BBB) shell. Endocrine anatomy was represented by a bilobed thyroid with an isthmus and four parathyroid surrogates. A right-sided bulky oropharyngeal-primary plus nodal tumour surrogate was generated as the union of two ellipsoidal volumes with exclusion of airway, spinal cord, and brainstem voxels. The gross tumour volume (GTV) was defined by this tumour mask, while the planning target volume (PTV) was created by geometric expansion using larger primary and nodal ellipsoids, constrained to the body mask and excluding the airway. A hypoxic subvolume was defined as nested ellipsoids within the GTV to support oxygen-dependent biological modulation. Figure~\ref{fig:voxel_phantom_anatomy} illustrates the final anatomical layout, and Figure~\ref{fig:voxel_phantom_density} shows the corresponding heterogeneous density map.

\subsection{Heterogeneous TOPAS ImageCube material phantom}
The voxelized anatomy was converted into a TOPAS \texttt{TsImageCube} phantom for Monte Carlo transport. Each voxel was assigned an explicit material tag mapped to a Geant4 base material plus a bulk density (Table~\ref{tab:voxel_materials}). Sixteen material classes were used: air, generic soft tissue, brain, BBB, brainstem, spinal cord, skull cortical bone, mandible/maxilla bone, vertebral bone, parotid, submandibular gland, thyroid, parathyroid, arterial blood, venous blood, and tumour. Representative bulk densities ranged from $0.0012$~g/cm$^3$ for air to $1.85$~g/cm$^3$ for skull cortical bone. Material tags were written to a binary ImageCube file, and a separate TOPAS material include was generated using \texttt{BaseMaterial} plus explicit \texttt{Density} assignments for each class. This produced a dose-calculation phantom that preserved both anatomical compartmentalization and tissue-density heterogeneity while remaining fully synthetic and reproducible. The materialized phantom used for Monte Carlo dose calculation is shown in Figure~\ref{fig:topas_material_phantom}.

\subsection{Direct SFRT lattice plan generation on the voxelized phantom}
A direct lattice-boost SFRT plan was constructed on the heterogeneous phantom. Candidate lattice vertices were required to lie fully within the GTV, with each candidate tested by enforcing that a spherical support region corresponding to the planned beamlet radius could be inscribed within the GTV. Candidate locations were sampled on a nominal $18 \times 20 \times 18$~mm lattice in the left-right, superior-inferior, and anterior-posterior directions about the target centroid. The accepted direct plan contained four valid lattice vertices. For delivery, the plan consisted of one broad anterior-posterior (AP) base field and a set of orthogonal lattice boosts delivered from AP, left lateral, and right lateral directions. The AP base-field radius was derived from the $99^\mathrm{th}$ percentile projected PTV radius plus a 6~mm geometric margin, and the lateral aperture size was derived analogously from the sagittal projection. Beam source planes were placed 5~mm beyond the phantom extent, yielding an AP source just anterior to the phantom and two lateral sources just outside the left and right surfaces. The base field received 95\% of the total histories, with the remaining 5\% distributed equally across the orthogonal lattice beamlets. All raw Monte Carlo dose distributions were finally normalized so that physical PTV $D_{95}$ matched a 6~Gy prescription. The treatment-plan geometry overlaid on the explicit arterial and venous network is shown in Figure~\ref{fig:treatment_plan_with_vessels}. Source-model and transport parameters are summarized in Table~\ref{tab:planning_transport_parameters}.

\subsection{Monte Carlo dose calculation}
Dose was calculated in TOPAS using a direct-photon polyenergetic surrogate of a nominal 6~MV clinical beam. The beam energy spectrum was sampled from a CSV-defined discrete 6~MV representative photon spectrum and emitted as a flat elliptical gamma beam with no additional angular spread. Transport was performed using the \texttt{em\_opt4\_only} modular Geant4 electromagnetic physics profile with production cuts of 0.01~mm for photons, electrons, and positrons. The detailed head-and-neck direct-plan simulations used $10^6$ primary histories, 8 threads, and a fixed seed of 33. Monte Carlo output was read back into Python as a voxelized physical dose grid and scaled so that the physical PTV $D_{95}$ equaled the nominal prescription dose. No patient-specific CT or clinical treatment-planning-system optimizer was used in this branch; instead, the phantom and source geometry were generated fully within the computational workflow to allow controlled testing of heterogeneous anatomy and explicit vascular structure.

\subsection{Biology-aware reinterpretation using explicit vascular sink fields}
The previously calibrated multiscale biology model was then transferred to the voxelized phantom without retuning of the locked transport parameters. Physical dose was first converted to a local linear-quadratic survival field using $\alpha=0.03$~Gy$^{-1}$ and $\beta=0.003$~Gy$^{-2}$. Two signaling species were then propagated using the coupled reaction-diffusion model: a short-range ROS-like channel and a longer-range cytokine-like channel. For the voxelized phantom branch, tissue-state tensors were constructed directly from anatomy. The cytokine emission multiplier was set to 2.0 within the GTV. Within the hypoxic subvolume, ROS emission was reduced to 0.12 of baseline and cytokine emission was increased by a factor of 2.7. Most importantly, the explicit vascular network embedded in the phantom was used directly as the spatial sink field. Arterial voxels were assigned uptake coefficients of 0.05 for ROS and 0.70 for cytokines, whereas venous voxels were assigned 0.05 for ROS and 0.90 for cytokines, thereby representing stronger cytokine washout in the venous compartment. Temporal integration was performed for 400 time steps with $\Delta t = 0.12$ using the locked transport settings $D_{\mathrm{cyto}} = 1.2$, $\lambda_{\mathrm{cyto}} = 0.001$, $\gamma = 0.35$, and scaling factor $s = 0.0029365813$, together with $D_{\mathrm{ROS}} = 0.8$, $\lambda_{\mathrm{ROS}} = 0.2$, $E_{\max,\mathrm{ROS}} = 1.5$, $E_{\max,\mathrm{cyto}} = 0.8$, $w_{\mathrm{ROS}} = 0.4$, $w_{\mathrm{cyto}} = 0.4$, and $w_{\mathrm{immune}} = 0.2$. The resulting final survival field was converted into an LQ-equivalent effective dose, $D_{\mathrm{eff}}$, by inversion of the linear-quadratic relation. The spatial cytokine sink field derived from the explicit arteries and veins is shown in Figure~\ref{fig:vascular_sink_field}. Biology-model parameters used in this voxelized branch are listed in Table~\ref{tab:bioaware_parameters}.

\subsection{Exploratory biology-guided lattice candidate optimization}
To begin closing the loop between planning and biological response, an exploratory biology-guided placement workflow was implemented on the voxelized phantom. The phase-15 direct plan was treated as the baseline candidate. Additional candidate lattice plans were then generated by sampling valid 8~mm-radius lattice centers throughout the GTV on a 6~mm candidate grid, enforcing a minimum inter-spot spacing of 18~mm and hard serial-organ avoidance relative to the spinal cord and brainstem. Candidate spots were scored using a feedback heuristic that combined tumour under-coverage in effective-dose space, hypoxia-aware boosting, proximity to vascular sink support, prior selection history, and distance-based penalties from organs at risk. Plans were then evaluated end-to-end with both physical and biology-aware metrics. A hard acceptance criterion required effective spinal cord $D_{2} \leq 45$~Gy and effective brainstem $D_{2} \leq 10$~Gy. Among plans that satisfied these safety constraints, the ranking objective favored high tumour effective dose and penalized excess OAR burden according to
\[
\mathcal{J} = 2.5\,D_{95,\mathrm{GTV}}^{(\mathrm{eff})} + 2.0\,D_{95,\mathrm{PTV}}^{(\mathrm{eff})} + 0.5\,D_{50,\mathrm{GTV}}^{(\mathrm{eff})} - \sum_j w_j \max\!\left(0, M_j^{(\mathrm{eff})} - T_j\right),
\]
where $M_j^{(\mathrm{eff})}$ denotes the relevant effective-dose metric for OAR $j$, $T_j$ is its acceptance threshold, and $w_j$ is a structure-specific penalty weight. The exploratory optimization stage therefore used the biology model as the primary decision variable for plan ranking rather than relying on physical DVHs alone. The optimization settings used so far are summarized in Table~\ref{tab:bioopt_parameters}.

\subsection{Dosimetric and biology-aware evaluation}
For each voxelized-phantom plan, structure-based metrics were computed from both the physical dose grid and the LQ-equivalent effective-dose grid. Target metrics included mean dose, $D_{98}$, $D_{95}$, $D_{50}$, $D_{2}$, $V_{95}$, $V_{100}$, and target coverage ratio. OAR metrics included mean dose, $D_{2}$, and selected $V_x$ values depending on structure type. Differential dose-volume data were converted into cumulative DVHs for the PTV, GTV, spinal cord, brainstem, parotids, mandible, thyroid, parathyroids, brain, and BBB. For biology-aware comparisons, the same endpoint definitions were retained but were applied to $D_{\mathrm{eff}}$ rather than physical dose. In addition, full survival fields were retained to support future assessment of cumulative target coverage, preservation of peak-valley structure, hotspot control, and inter-fraction consistency under fractionated plan sequences.
